% minilecture04.tex: Tutorial 4 mini-lecture: Surfaces & Interiors
% Planetary Systems, Kapteyn Institute, University of Groningen
% A 5-10 minute recap, presented by the TAs at the start of Tutorial 4,
% of the Lecture 7-8 concepts that Worksheet 4 trains.

\documentclass[aspectratio=169, 11pt]{beamer}

% ── Theme and macros (shared with the lecture decks) ────────
\usepackage{../../slides/common/beamerthemeIPS}
% macros.tex: Shared math macros for IPS lecture slides
% Mirrors definitions in book/_config.yml for consistency
% Uses \providecommand to avoid conflicts with unicode-math

% ── Derivatives ─────────────────────────────────────────────
\providecommand{\dv}[2]{\frac{\mathrm{d} #1}{\mathrm{d} #2}}
\providecommand{\dd}{}\renewcommand{\dd}{\mathrm{d}}
\providecommand{\pdv}[2]{\frac{\partial #1}{\partial #2}}

% ── Solar quantities ────────────────────────────────────────
\newcommand{\Msun}{M_{\odot}}
\newcommand{\Rsun}{R_{\odot}}
\newcommand{\Lsun}{L_{\odot}}

% ── Planetary quantities ────────────────────────────────────
\newcommand{\Mearth}{M_{\oplus}}
\newcommand{\Rearth}{R_{\oplus}}
\newcommand{\Mjup}{M_{\mathrm{J}}}
\newcommand{\Rjup}{R_{\mathrm{J}}}

% ── Constants ───────────────────────────────────────────────
\newcommand{\kB}{k_{\mathrm{B}}}

% ── Media links ─────────────────────────────────────────────
% Clickable button that opens the animated or video version of a
% figure, e.g. \mediabutton{https://...gif}{Play animation}
\providecommand{\mediabutton}[2]{\href{#1}{\beamergotobutton{#2}}}

\graphicspath{{figures/}{../../slides/lecture08/figures/}{../../slides/lecture07/figures/}{../../slides/common/}}

% ── Metadata ────────────────────────────────────────────────
\title[T4 mini-lecture: Surfaces \& Interiors]{Tutorial 4 Mini-Lecture:\\Surfaces \& Interiors}
\subtitle{The concepts behind Worksheet 4 \textbullet\ Lectures 7--8}
\author{}
\institute{Kapteyn Astronomical Institute\\University of Groningen}
\date{2026}
\titlebgcredit{Earth's interior cross-section. Credit: NASA/JPL-Caltech}
\titlebgopacity{0.55}

% ════════════════════════════════════════════════════════════
\begin{document}

% ── Title slide ─────────────────────────────────────────────
\begin{frame}[plain,noframenumbering]
  \titlepage
\end{frame}

% ── Roadmap ─────────────────────────────────────────────────
\begin{frame}{What Worksheet 4 trains}
  Five problems, all built on Lectures 7--8. The physics you need, per problem:
  \vskip4pt\relax
  \begin{enumerate}
    \item \textbf{Impact energetics on Mars}: impactor mass and energy, gravity-regime crater scaling, the simple-to-complex transition.
    \item \textbf{Reading ages from craters}: the crater chronology $N(T)$, its linear and exponential branches, the printed chronology figure.
    \item \textbf{Moment of inertia and the inside of Mercury}: the two-layer model, layer densities, what $C/MR^2$ measures.
    \item \textbf{Central pressure}: hydrostatic equilibrium, the uniform-density pressure profile, why it is a lower bound.
    \item \textbf{Seismology}: shadow zones, what liquid layers do to S and P waves, the Adams-Williamson assumptions.
  \end{enumerate}
  \vskip2pt\relax
  \keyresult{Every problem has the shape and level of one exam question. All parts are analytical; no computer is needed.}
\end{frame}

% ── P1 ──────────────────────────────────────────────────────
\begin{frame}{Impact energetics on Mars}
  \small
  \begin{columns}[T]
    \column{0.52\textwidth}
    The gravity-regime scaling of Lecture 7:
    \[
      \boxed{D \approx \left(\frac{E_k}{\rho_t\,g}\right)^{1/4}}
      \qquad D \propto L^{3/4}
    \]
    \begin{itemize}
      \item $E_k = \tfrac{1}{2}mv^2$ with $m = \tfrac{4}{3}\pi r^3 \rho$: the whole chain runs in SI units.
      \item The fourth root compresses hard: check how many orders of magnitude in energy correspond to one order of magnitude in diameter.
    \end{itemize}

    \column{0.44\textwidth}
    \begin{block}{Simple or complex?}
      The transition diameter scales as $D_t \propto 1/g$: stronger gravity collapses smaller cavities. Locate the martian transition from the lunar one, then classify your crater against it.
    \end{block}
  \end{columns}
  \vskip2pt\relax
  \keyresult{One scaling law connects the impactor you cannot see to the crater you can measure.}
\end{frame}

% ── P2 ──────────────────────────────────────────────────────
\begin{frame}{Reading ages from craters}
  \small
  \begin{columns}[T]
    \column{0.52\textwidth}
    The lunar chronology of Lecture 7:
    \[
      \boxed{N(\geq 1\,\mathrm{km}) = a\left(e^{\lambda T} - 1\right) + b\,T}
    \]
    \begin{itemize}
      \item Young surfaces lie on the linear branch: steady flux, $N \approx bT$, age is a division.
      \item Old surfaces lie on the exponential branch of the early bombardment: $N$ grows exponentially with $T$.
    \end{itemize}

    \column{0.44\textwidth}
    \begin{block}{Iterate, do not invert}
      The full $N(T)$ has no closed-form inverse. On the steep branch, neglect the small term, solve, then correct once with the term you dropped. On the shallow branch, check the dropped term is small instead.
    \end{block}
  \end{columns}
  \vskip2pt\relax
  \keyresult{The same measurement error means different things on different branches: check which branch you are on before you trust a crater age.}
\end{frame}

% ── P3 ──────────────────────────────────────────────────────
\begin{frame}{Moment of inertia and the inside of Mercury}
  \small
  \begin{columns}[T]
    \column{0.52\textwidth}
    The two-layer model of Lecture 8:
    \[
      \boxed{\frac{C}{MR^2} = \frac{2}{5}\,\frac{(f-1)\,x^5 + 1}{(f-1)\,x^3 + 1}}
      \quad f = \frac{\rho_c}{\rho_m}
    \]
    \begin{itemize}
      \item A uniform sphere gives $0.400$; every planet measures below it, and the deficit measures central condensation.
      \item Mass balance fixes the layer densities: core mass in core volume, the rest in the mantle shell.
    \end{itemize}

    \column{0.44\textwidth}
    \begin{block}{Small cores hide well}
      A core of radius fraction $x$ fills only $x^3$ of the volume. Before judging any claim from $C/MR^2$, compute how far a plausible core can move the number at all.
    \end{block}
  \end{columns}
  \vskip2pt\relax
  \keyresult{$C/MR^2$ is measured from spin and gravity without drilling a hole; it is the single most informative number about a planetary interior.}
\end{frame}

% ── P4 ──────────────────────────────────────────────────────
\begin{frame}{Central pressure}
  \small
  \begin{columns}[T]
    \column{0.52\textwidth}
    Hydrostatic equilibrium at uniform density:
    \[
      \boxed{P_c = \frac{2\pi}{3}\,G\bar\rho^2 R^2 = \frac{3GM^2}{8\pi R^4}}
    \]
    \begin{itemize}
      \item Derive it once: integrate $\dd P/\dd r = -Gm(r)\rho/r^2$ from the centre to the surface.
      \item The two forms must agree; use the second as a cross-check on the first.
    \end{itemize}

    \column{0.44\textwidth}
    \begin{block}{What changes if density is not uniform?}
      Compare the enclosed-mass term at fixed $r$ for a centrally concentrated profile against a uniform one. State only which term grows, not the direction of the final inequality.
    \end{block}
  \end{columns}
  \vskip2pt\relax
  \keyresult{The central pressure of a uniform body depends only on its mass and radius; the strong $M^2/R^4$ scaling separates the planets by orders of magnitude.}
\end{frame}

% ── P5 ──────────────────────────────────────────────────────
\begin{frame}{Seismology: listening to interiors}
  \small
  \begin{columns}[T]
    \column{0.52\textwidth}
    The wave speeds of Lecture 8:
    \[
      \boxed{v_S = \sqrt{\frac{\mu}{\rho}}
      \qquad
      v_P = \sqrt{\frac{K + \tfrac{4}{3}\mu}{\rho}}}
    \]
    \begin{itemize}
      \item A liquid has no shear modulus: $\mu = 0$ stops S waves entirely.
      \item P waves pass through liquids, but a speed drop at a boundary bends the rays and opens a shadow ring.
    \end{itemize}

    \column{0.44\textwidth}
    \begin{block}{Two shadows, two mechanisms}
      One shadow comes from absorption, the other from refraction. For each observation, state which mechanism makes it and which interior property it therefore proves.
    \end{block}
  \end{columns}
  \vskip2pt\relax
  \keyresult{What is missing from a seismogram tells as much as what arrives: shadow zones map liquid layers on any planet with a seismometer.}
\end{frame}

% ── Closer ──────────────────────────────────────────────────
\begin{frame}{Working the worksheet}
  \begin{itemize}
    \item \textbf{Show your algebra.} The exam asks for the steps, not only the number.
    \item Carry \textbf{4 significant figures} through intermediate steps; round at the end.
    \item Use the \textbf{checkpoints} ("show that \ldots") to verify your work and to keep moving if a step resists.
    \item Qualitative parts want \textbf{2 to 3 sentences} of physics, not an essay.
    \item Problem 2 reads the chronology figure printed on the sheet; Problem 5 is entirely qualitative and trains the exam skill of arguing from mechanisms.
  \end{itemize}
  \vskip8pt\relax
  \keyresult{Worksheet and full solutions: \texttt{ips.formingworlds.space/worksheets.html}. We discuss questions, not presentations of the solutions.}
\end{frame}

\end{document}
